% GENERATED FILE. Edit canonical lessons, then run npm run generate:books.
% canonical-source-hash: fnv1a64:960713f8caf5b17d
% canonical-lessons: LA-C371-rostra, LA-C371-contio, LA-C371-pronuntio, LA-C371-recito, LA-C371-temere

\chapter{At the Speakers' Platform}
\label{ch:la-at-the-speakers-platform}
\hlchaptermodality{371}

\begin{chapteropening}
\textbf{By the end of this chapter:} \emph{I can talk about the speakers' platform and a public meeting, and say that someone announced, read aloud or spoke rashly.}

Complete the last lesson of chapter 371: I can talk about the speakers' platform and a public meeting, and say that someone announced, read aloud or spoke rashly.
\end{chapteropening}

\section[rōstra, rōstrōrum]{rōstra, rōstrōrum — the speakers' platform (in the Forum)}

\label{lesson:LA-C371-rostra}



\begin{quote}
Before the new one: say the Latin for lightness, then the Latin for rashness.
\end{quote}

\subsection*{You'll want to know: rōstra, rōstrōrum}
\textbf{rōstra, rōstrōrum} — \textquotedblleft{}the speakers' platform (in the Forum)\textquotedblright{}.

\begin{quote}
Senex in rōstra ascendit.
\end{quote}

\begin{quote}
\emph{The old man went up onto the speakers' platform.}
\end{quote}

\emph{rōstra, rōstrōrum} n. pl. The platform in the Forum was decked with the bronze beaks (\emph{rōstra}) of captured warships. English \emph{rostrum}.

\begin{cousinweb}
\emph{rōstra} is the plural of \emph{rōstrum} \textquotedblleft{}a beak, a ship's ram\textquotedblright{}, which is built on \emph{rōdō} \textquotedblleft{}to gnaw\textquotedblright{} (Lewis \& Short). English \emph{rostrum} comes from it.
\end{cousinweb}

\begin{grammarlens}[title={What you've built}]
One more word for telling, explaining and describing at B1.
\end{grammarlens}

\subsection*{Your turn}
\begin{itemize}
\raggedright
  \item \emph{Say it:} \emph{rōstra, rōstrōrum}
  \item \emph{Say it:} \emph{rōstra, rōstrōrum}, once more
  \item \emph{Say it:} a sentence of your own with \emph{rōstra, rōstrōrum}
  \item \emph{Recall:} say the Latin for lightness, then the Latin for rashness, then say \emph{rōstra, rōstrōrum} again
\end{itemize}

\begin{tcolorbox}[breakable,colback=teal!4,colframe=teal!35!black,title={Before you move on}]
What does rōstra, rōstrōrum mean? (\textquotedblleft{}The speakers' platform (in the Forum)\textquotedblright{}.) Say it once more.
\end{tcolorbox}
\section[cōntiō, cōntiōnis]{cōntiō, cōntiōnis — a public meeting; a speech made to it}

\label{lesson:LA-C371-contio}



\begin{quote}
Before the new one: say the Latin for rashness, then the Latin for the speakers' platform.
\end{quote}

\subsection*{You'll want to know: cōntiō, cōntiōnis}
\textbf{cōntiō, cōntiōnis} — \textquotedblleft{}a public meeting; a speech made to it\textquotedblright{}.

\begin{quote}
Cōnsul cōntiōnem habuit.
\end{quote}

\begin{quote}
\emph{The consul held a public meeting.}
\end{quote}

\emph{cōntiō, cōntiōnis} f. A meeting of the people called to hear a magistrate, where no vote was taken; also the speech made there.

\begin{cousinweb}
\emph{cōntiō} is contracted from an old form \emph{coventiō} \textquotedblleft{}a coming together\textquotedblright{} (Lewis \& Short), built from \emph{con-} and \emph{veniō} \textquotedblleft{}to come\textquotedblright{}.
\end{cousinweb}

\begin{grammarlens}[title={What you've built}]
One more word for telling, explaining and describing at B1.
\end{grammarlens}

\subsection*{Your turn}
\begin{itemize}
\raggedright
  \item \emph{Say it:} \emph{cōntiō, cōntiōnis}
  \item \emph{Say it:} \emph{cōntiō, cōntiōnis}, once more
  \item \emph{Say it:} a sentence of your own with \emph{cōntiō, cōntiōnis}
  \item \emph{Recall:} say the Latin for rashness, then the Latin for the speakers' platform, then say \emph{cōntiō, cōntiōnis} again
\end{itemize}

\begin{tcolorbox}[breakable,colback=teal!4,colframe=teal!35!black,title={Before you move on}]
What does cōntiō, cōntiōnis mean? (\textquotedblleft{}A public meeting; a speech made to it\textquotedblright{}.) Say it once more.
\end{tcolorbox}
\section[prōnūntiō, prōnūntiāre]{prōnūntiō, prōnūntiāre — to proclaim, to announce; to deliver (a speech)}

\label{lesson:LA-C371-pronuntio}



\begin{quote}
Before the new one: say the Latin for the speakers' platform, then the Latin for a public meeting.
\end{quote}

\subsection*{You'll want to know: prōnūntiō, prōnūntiāre}
\textbf{prōnūntiō, prōnūntiāre} — \textquotedblleft{}to proclaim, to announce; to deliver (a speech)\textquotedblright{}.

\begin{quote}
Praecō sententiam iūdicum prōnūntiāvit.
\end{quote}

\begin{quote}
\emph{The herald announced the judges' verdict.}
\end{quote}

\emph{prōnūntiō, prōnūntiāre, prōnūntiāvī, prōnūntiātum}. Built on \emph{nūntiō}, to announce. It is also a speaker's delivery. English \emph{pronounce}.

\begin{cousinweb}
\emph{prōnūntiō} is built from \emph{prō} \textquotedblleft{}forth\textquotedblright{} and \emph{nūntiō} \textquotedblleft{}to announce\textquotedblright{} (Lewis \& Short). English \emph{pronounce} comes from it through French.
\end{cousinweb}

\begin{grammarlens}[title={What you've built}]
One more word for telling, explaining and describing at B1.
\end{grammarlens}

\subsection*{Your turn}
\begin{itemize}
\raggedright
  \item \emph{Say it:} \emph{prōnūntiō, prōnūntiāre}
  \item \emph{Say it:} \emph{prōnūntiō, prōnūntiāre}, once more
  \item \emph{Say it:} a sentence of your own with \emph{prōnūntiō, prōnūntiāre}
  \item \emph{Recall:} say the Latin for the speakers' platform, then the Latin for a public meeting, then say \emph{prōnūntiō, prōnūntiāre} again
\end{itemize}

\begin{tcolorbox}[breakable,colback=teal!4,colframe=teal!35!black,title={Before you move on}]
What does prōnūntiō, prōnūntiāre mean? (\textquotedblleft{}To proclaim, to announce; to deliver (a speech)\textquotedblright{}.) Say it once more.
\end{tcolorbox}
\section[recitō, recitāre]{recitō, recitāre — to read aloud, to recite}

\label{lesson:LA-C371-recito}



\begin{quote}
Before the new one: say the Latin for a public meeting, then the Latin for to proclaim.
\end{quote}

\subsection*{You'll want to know: recitō, recitāre}
\textbf{recitō, recitāre} — \textquotedblleft{}to read aloud, to recite\textquotedblright{}.

\begin{quote}
Poēta carmina sua amīcīs recitāvit.
\end{quote}

\begin{quote}
\emph{The poet read his poems aloud to his friends.}
\end{quote}

\emph{recitō, recitāre, recitāvī, recitātum}. At Rome authors read new work aloud to invited listeners. English \emph{recite}.

\begin{cousinweb}
\emph{recitō} is built from \emph{re-} and \emph{citō} \textquotedblleft{}to call, to summon\textquotedblright{} (EtymDB; Lewis \& Short): to call words out again. English \emph{recite} comes from it.
\end{cousinweb}

\begin{grammarlens}[title={What you've built}]
One more word for telling, explaining and describing at B1.
\end{grammarlens}

\subsection*{Your turn}
\begin{itemize}
\raggedright
  \item \emph{Say it:} \emph{recitō, recitāre}
  \item \emph{Say it:} \emph{recitō, recitāre}, once more
  \item \emph{Say it:} a sentence of your own with \emph{recitō, recitāre}
  \item \emph{Recall:} say the Latin for a public meeting, then the Latin for to proclaim, then say \emph{recitō, recitāre} again
\end{itemize}

\begin{tcolorbox}[breakable,colback=teal!4,colframe=teal!35!black,title={Before you move on}]
What does recitō, recitāre mean? (\textquotedblleft{}To read aloud, to recite\textquotedblright{}.) Say it once more.
\end{tcolorbox}
\section[temerē]{temerē — rashly, blindly; at random}

\label{lesson:LA-C371-temere}



\begin{quote}
Before the new one: say the Latin for to proclaim, then the Latin for to read aloud.
\end{quote}

\subsection*{You'll want to know: temerē}
\textbf{temerē} — \textquotedblleft{}rashly, blindly; at random\textquotedblright{}.

\begin{quote}
Nihil temerē faciēbat.
\end{quote}

\begin{quote}
\emph{He did nothing rashly.}
\end{quote}

An adverb; its noun is \emph{temeritās}, rashness. \emph{Nōn temerē}, not without good reason.

\begin{cousinweb}
\emph{temerē} is inherited from Proto-Italic \emph{*temezi} \textquotedblleft{}in the dark\textquotedblright{} (EtymDB), kin to \emph{tenebrae} \textquotedblleft{}darkness\textquotedblright{} (Lewis \& Short): to act blindly.
\end{cousinweb}

\begin{grammarlens}[title={What you've built}]
One more word for telling, explaining and describing at B1.
\end{grammarlens}

\subsection*{Your turn}
\begin{itemize}
\raggedright
  \item \emph{Say it:} \emph{temerē}
  \item \emph{Say it:} \emph{temerē}, once more
  \item \emph{Say it:} a sentence of your own with \emph{temerē}
  \item \emph{Recall:} say the Latin for to proclaim, then the Latin for to read aloud, then say \emph{temerē} again
\end{itemize}

\begin{tcolorbox}[breakable,colback=teal!4,colframe=teal!35!black,title={Before you move on}]
What does temerē mean? (\textquotedblleft{}Rashly, blindly; at random\textquotedblright{}.) Say it once more.
\end{tcolorbox}
